\section{本章小结}
\label{sec:ls-representation-summary}

本章把对象与记录对象的数分开。线性空间规定哪些对象可以相加和缩放，子空间是其中对线性组合封闭的部分。子空间可以求交、求和；交为零时和是直和，分解唯一。商空间则把相差某个子空间中向量的对象视为相同。

基同时满足张成与线性无关，所以每个向量有唯一的坐标。基的向量个数与选取无关，这个数就是维数；“先取已有的基，再补新方向”给出和空间、补空间与商空间的维数。换基改变坐标，不改变向量，新旧坐标由可逆的换基矩阵换算。基接近退化时，坐标可能很大，不能用来判断对象的大小。

线性映射保持线性组合。选定两端的基，它记成矩阵，叠加、复合与求逆分别对应矩阵的加法、乘法与求逆。核记录丢失的输入，像记录能到达的输出，秩零度定理说明二者的维数之和等于输入维数；由此可以判断单射、满射与可逆。同一个映射在不同基下的矩阵相差两侧的可逆因子：两端各自换基得到相抵，秩是它的完整不变量；算子两端同时换基得到相似，秩、行列式与迹相同仍不足以保证相似。

对偶空间由线性泛函组成，对偶基就是坐标读取。对偶映射把输出端的读取传回输入端，在对偶基下记录为转置，由此又得到行秩等于列秩。

这些结构还没有规定向量的大小或夹角。下一章引入范数和内积，讨论收敛、正交分解与最佳近似。
